\documentclass[11pt]{article}

\usepackage[margin=1in]{geometry}
\usepackage{amsmath, amssymb}
\usepackage{graphicx}
\usepackage{float}
\usepackage[hidelinks]{hyperref}
% Long file:line paths in \texttt would otherwise overflow the margin.
\usepackage[htt]{hyphenat}
\sloppy

% Figures live in the git-ignored output tree, referenced relative to this file.
\graphicspath{{../../output/}}

\newcommand{\Ainv}{A^{-1}}

\title{Paper story: experiments}
%\author{Marius }
\date{\today}

\begin{document}
\maketitle

\noindent This is a progress report: the story of the paper told through the figures we have
produced so far. It is organised by PDE --- steady-state Poisson first, then the time-dependent
Allen--Cahn stepper --- and closes with what comes next (Stokes). \textbf{I am sorry for the language, but making Claude make this is so much faster.}

\section{Poisson}

\subsection{Setup: the discrete problem, and where it lives in the code}
\label{sec:poisson-setup}

Everything below is what the code actually does, checked against the source rather than
reconstructed from memory. File references are \texttt{path:line} at the time of writing.
The reference configuration is the one both Poisson figures come from
(\texttt{experiments/poisson/sweep\_blend/pls/sweep.sbatch}).

\paragraph{Continuous problem and finite element space.}
We solve $-\Delta u = f$ on $\Omega = (0,1)^2$ with homogeneous Dirichlet boundary conditions.
The discretisation is the space of \textbf{bilinear quadrilateral ($Q_1$) elements} on a
structured tensor-product grid, built by \texttt{structured\_quad\_mesh}
(\texttt{tensorpils/meshing.py:24}) as a genuine \texttt{tensormesh.Mesh} with a single
\texttt{"quad"} cell block and row-major node numbering $k = i\,n_x + j$. At the default
$64\times 64$ nodes this gives $N = 4096$ degrees of freedom, $63^2 = 3969$ cells, node spacing
$h = 1/63$, and $252$ boundary nodes. The stiffness $A$ and mass $M$ are assembled once, at
construction, by TensorMesh's \texttt{LaplaceElementAssembler} and \texttt{MassElementAssembler}
(\texttt{tensorpils/physics.py:64--65}) using Gauss quadrature of order $4$ per axis
(\texttt{physics.py:59}), which integrates the $Q_1$ stiffness and mass exactly on rectangular
cells.

\paragraph{How the boundary condition is imposed.}
The Dirichlet condition is \textbf{not} statically condensed. Instead boundary nodes are
\emph{masked}: the prediction is projected to zero on the boundary, the residual is formed, and
the residual is zeroed on the boundary again (\texttt{apply\_zero\_boundary},
\texttt{physics.py:32}; \texttt{PoissonProblem.residual}, \texttt{physics.py:100--108}),
\begin{equation*}
  r(u,f) \;=\; \Pi_0\big(A\,\Pi_0 u \;-\; M f\big),
\end{equation*}
where $\Pi_0$ zeroes boundary entries. The operator that effectively acts is therefore the
\emph{interior block} $A_{II}$ of size $3844 \times 3844$, and this is also the block the spectral
preconditioner is built from. Note that the right-hand side is the \textbf{consistent finite
element load} $b = Mf$ (\texttt{load\_vector}, \texttt{physics.py:83}), not the pointwise nodal
$f$ --- which matters, because $\tfrac12\|Au-b\|^2$ and $\tfrac12\|Au - f\|^2$ are different
problems.

\paragraph{Distribution over right-hand sides.}
Sources are random truncated sine series (\texttt{tensorpils/data.py:74--78}, wrapping
TensorMesh's \texttt{PoissonMultiFrequency}). We draw $a_{ij} \sim \mathrm{Unif}[-1,1]$
i.i.d.\ for $i,j = 1,\dots,K$ and set
\begin{align*}
  f(x,y) &\;=\; \frac{\pi}{K^2}\sum_{i,j=1}^{K} a_{ij}\,(i^2+j^2)^{1/2}\,\sin(\pi i x)\sin(\pi j y), \\
  u^\star(x,y) &\;=\; \frac{1}{\pi K^2}\sum_{i,j=1}^{K} a_{ij}\,(i^2+j^2)^{-1/2}\,\sin(\pi i x)\sin(\pi j y),
\end{align*}
so that $-\Delta u^\star = f$ exactly. All experiments here use $\mathbf{K=4}$, i.e.\ $16$ active
modes with maximum frequency $4$ --- very comfortably resolved on the $64^2$ grid. Two points are
worth stating explicitly because they are easy to get backwards:
\begin{itemize}
  \item The source is \textbf{amplified} at high frequency (weight $|k| = (i^2+j^2)^{1/2}$) while
        the solution \emph{decays} like $|k|^{-1}$. The sources are rough, the solutions smooth.
  \item TensorMesh's own docstring for \texttt{source\_term} advertises the exponent as
        $(i^2+j^2)^{r}$ with $r=-0.5$, i.e.\ the \emph{opposite} tilt. The code computes
        $(i^2+j^2)^{-r}$ (\texttt{tensormesh/dataset/equation/poisson.py:81}). The code is
        internally consistent --- I verified numerically that $-\Delta u^\star = f$ holds for the
        pair as implemented --- so this is a documentation bug upstream, not a data bug. Do not
        quote the docstring formula in the paper.
\end{itemize}
The reference $u^\star$ is the \textbf{analytical} solution sampled at the nodes, \emph{not} the
Galerkin solution $A^{-1}b$. The reported relative $L^2$ therefore carries a floor at the $Q_1$
discretisation error; at $K=4$ on $64^2$ that floor is orders of magnitude below the errors
plotted, but it is not zero, and it will matter if we ever push these curves much lower.

\paragraph{Splits.} A single pool of samples is generated once from \texttt{seed=42} and sliced
sequentially into $n_{\text{train}} = 1024$, $n_{\text{val}} = 128$, $n_{\text{test}} = 256$
(\texttt{data.py:99--121}); all three splits share one mesh and one assembled operator.

\paragraph{The loss that is actually optimised.}
For \texttt{--loss pls} (\texttt{PreconditionedLSLoss}, \texttt{tensorpils/losses.py:159}),
\begin{equation*}
  \mathcal{L}(\theta) \;=\; \frac{1}{B}\sum_{b=1}^{B} \frac12 \big\| P\,r(u_\theta(f_b), f_b) \big\|_2^2,
\end{equation*}
i.e.\ \textbf{summed over nodes, averaged over the batch}. Because the residual hard-projects the
boundary, the PLS loss ignores \texttt{--lambda\_bc} and \texttt{--bc\_mode} entirely; there is no
boundary penalty term in these runs.

\paragraph{The multigrid V-cycle.}
\texttt{GeometricMultigrid} (\texttt{tensorpils/preconditioners/multigrid.py}) supplies
$P \approx \Ainv$ as \textbf{exactly one V-cycle} per loss evaluation, with the CLI defaults:
\begin{itemize}
  \item \textbf{Hierarchy:} $4$ levels, coarsening $(n+1)//2$ per axis (floored at $3$). At $64^2$
        the levels are $64^2 \to 32^2 \to 16^2 \to 8^2$ (verified by instantiation).
  \item \textbf{Level operators are re-discretised, not Galerkin.} Each level assembles its own
        stiffness on its own structured mesh via \texttt{LaplaceElementAssembler}
        (\texttt{multigrid.py:111--125}); we do \emph{not} form $R A P$. Dirichlet rows and columns
        are then zeroed with a unit diagonal, so a residual with zero boundary entries keeps a zero
        boundary through the whole cycle.
  \item \textbf{Smoother:} weighted Jacobi, $\omega = 2/3$, $2$ pre- and $2$ post-smoothing sweeps
        (\texttt{multigrid.py:148--153}). Symmetric pre/post counts make $P$ symmetric, which is
        what the $\tfrac12\|Pr\|^2$ form wants.
  \item \textbf{Transfer:} bilinear prolongation respecting the row-major node order, restriction
        $R = P^\top$ (\texttt{multigrid.py:44--48}).
  \item \textbf{Coarsest level:} solved exactly, by a dense inverse precomputed at construction
        (\texttt{multigrid.py:97}) --- at $8^2$ that is cheap.
  \item \textbf{Differentiability:} the cycle is a sequence of sparse mat--vecs and element-wise
        operations, so autograd flows through it and the gradient of $\tfrac12\|Pr\|^2$ is exact
        rather than a surrogate.
\end{itemize}

\paragraph{The spectral (blend) preconditioner.}
For the $t$-sweep, $P$ is built as an exact function of the interior stiffness,
$P = Q\,p(\Lambda)\,Q^\top$ with $p(\lambda) = (1-t) + t/\lambda$
(\texttt{spectral.py:29}), which is precisely $P_t = (1-t)I + t\Ainv$ on the interior block. The
eigendecomposition is always taken in \texttt{float64}. At $64^2$ this is the dense
\texttt{eigh} (\texttt{--precond\_method dense}, the default); at $128^2$/$256^2$ the dense solve is
infeasible and the identical operator is realised through the fast discrete sine transform
(\texttt{--precond\_method sine}). Measured on the $64^2$ interior block: $\kappa(A_{II}) = 8.04\times 10^{2}$
and $\kappa(H) = \kappa(A)^2 = 6.47\times 10^{5}$, which is where the numbers quoted in the next
two subsections come from.

\paragraph{Model, optimisation, and the reported metric.}
The neural operator is \texttt{neuralop}'s FNO (\texttt{tensorpils/models.py}) with
$16\times 16$ Fourier modes, width $64$, and $5$ layers. Training is Adam with cosine-annealed
learning rate $10^{-3} \to 10^{-4}$, $500$ epochs, batch size $32$, no weight decay. The metric in
every figure is the \textbf{mass-matrix (FEM) relative $L^2$},
$\|u_\theta - u^\star\|_M / \|u^\star\|_M$, computed per sample and then averaged over the split
(\texttt{trainer.py:367--373}) --- not a plain Euclidean norm on nodal values. At evaluation the
prediction's boundary is projected to zero for \emph{every} loss (\texttt{trainer.py:303}), so no
method is charged for a boundary the BC already fixes.

\begin{table}[H]
  \centering
  \begin{tabular}{ll}
    \hline
    Domain / space & $(0,1)^2$, bilinear $Q_1$ on a structured quad grid \\
    Grid & $64\times 64$ nodes, $N=4096$, $h=1/63$, $3844$ interior dofs \\
    Quadrature & Gauss order $4$ per axis (exact for $Q_1$ on rectangles) \\
    BC & homogeneous Dirichlet, imposed by masking (not condensation) \\
    Right-hand side & $K=4$ sine series, $a_{ij}\sim\mathrm{Unif}[-1,1]$, weight $(i^2+j^2)^{1/2}$ \\
    Reference & analytical $u^\star$ (not the Galerkin solution) \\
    Load vector & consistent, $b = Mf$ \\
    Splits / seed & $1024$ / $128$ / $256$, seed $42$ \\
    Model & FNO, modes $16^2$, width $64$, $5$ layers \\
    Optimiser & Adam, cosine $10^{-3}\!\to\!10^{-4}$, $500$ epochs, batch $32$ \\
    Multigrid & $4$ levels, weighted Jacobi $\omega=2/3$, $2{+}2$ sweeps, one V-cycle \\
    Metric & mass-matrix relative $L^2$, per sample, averaged \\
    Conditioning & $\kappa(A_{II}) = 8.04\times10^{2}$, $\kappa(H) = 6.47\times10^{5}$ \\
    \hline
  \end{tabular}
  \caption{Reference Poisson configuration. Source:
  \texttt{experiments/poisson/sweep\_blend/pls/sweep.sbatch} plus CLI defaults in
  \texttt{tensorpils/cli.py}.}
  \label{tab:poisson-setup}
\end{table}

\subsection{Least-squares loss and spectral preconditioning}

\begin{itemize}
  \item Physics-informed neural operators often \textbf{do not train}, and the culprit is
        \emph{spectral ill-conditioning} inherited from the unbounded PDE operator: the loss
        landscape is a long, narrow valley that first-order optimisers cannot descend.
  \item To demonstrate this cleanly we use the \textbf{least-squares residual loss} based on the stiffness matrix $A$ on a $Q1$ finite element space. We insert a
        parametrised preconditioner $P_t$,
        \begin{equation*}
          \mathcal{L}(\theta) \;=\; \tfrac12\,\big\| P_t\,(A\,u_\theta - b) \big\|_2^2,
          \qquad
          P_t \;=\; (1-t)\,I \;+\; t\,\Ainv,
        \end{equation*}
        so $P_0 = I$ (no preconditioning) and $P_1 = \Ainv$ (full preconditioning), with a convex
        interpolation in between. The interpolation is made tractable either through the discrete
        sine transform (exact on the regular grid) or a dense eigenvalue solve.
  \item The sweep over $t$ (Figure~\ref{fig:pls-overlay}) shows the two extremes: \textbf{without}
        preconditioning ($t=0$) training \textbf{stalls}; with \textbf{full} preconditioning
        ($t=1$) we \textbf{recover the gold standard of data-driven training}. The reason is the
        squaring: the least-squares Hessian conditioning scales like $h^{-4}$ --- already
        $\approx 6.5\times 10^{5}$ on a $64^2$ grid, which is horrible.
  \item In practice $\Ainv$ is never formed; a \textbf{multigrid} V-cycle supplies $P\approx\Ainv$
        at \emph{optimal complexity}. The multigrid-preconditioned run matches data-driven training
        (rightmost curve in Figure~\ref{fig:pls-overlay}). The whole story is summarised by plotting
        the final validation error \emph{against the condition number} (Figure~\ref{fig:pls-collapse}):
        the runs collapse onto a single monotone curve, error rising with conditioning.
\end{itemize}

\begin{figure}[H]
  \centering
  \includegraphics[width=0.92\linewidth]{poisson/sweep_blend/pls/sweep_overlay.png}
  \caption{Poisson, least-squares loss. Validation error vs.\ epoch as the preconditioner strength
  $t$ is swept from $0$ (identity) to $1$ ($\Ainv$), plus the multigrid V-cycle and the data-driven
  reference. Unpreconditioned training stalls; full / multigrid preconditioning recovers the
  data-driven gold standard.}
  \label{fig:pls-overlay}
\end{figure}

\begin{figure}[H]
  \centering
  \includegraphics[width=0.80\linewidth]{poisson/sweep_blend/pls/sweep_collapse.png}
  \caption{The same runs, final validation error against the condition number of the preconditioned
  operator. Error grows monotonically with conditioning --- the mechanism, made explicit.}
  \label{fig:pls-collapse}
\end{figure}

\subsection{Deep Ritz: the variational formulation is milder}

The story changes for the \textbf{variational (Deep Ritz) formulation}: its conditioning scales
like $h^{-2}$ rather than $h^{-4}$. This is mild enough that on a $64^2$ grid training
\textbf{succeeds without any preconditioning} (Figure~\ref{fig:deepritz-overlay}). Preconditioning
can still be added and yields a small further speed-up --- it reduces the conditioning from
$\approx 800$ down to $1$. (On finer grids the $h^{-2}$ growth does become a real problem; there is
a figure for that, but it is not important enough to include here.)

\begin{figure}[H]
  \centering
  \includegraphics[width=0.92\linewidth]{poisson/sweep_blend/deepritz/sweep_overlay.png}
  \caption{Poisson, Deep Ritz (variational) loss. Because the conditioning is only $h^{-2}$,
  training already succeeds unpreconditioned on $64^2$; preconditioning gives a modest speed-up.}
  \label{fig:deepritz-overlay}
\end{figure}

\subsection{The infinite-data limit}

Because we generate \textbf{no labelled data}, we are not bound to a fixed dataset. For
Figure~\ref{fig:data-scaling} we fix the optimisation budget (500 epochs' worth of forward passes
at a dataset of 1024) and grow the training-set size from $32$ up to $4096$. Crucially, we also add
the \textbf{infinite dataset} --- one that never revisits a sample --- as the dashed black line.

This is a genuine selling point of physics-informed operator learning: you want to maximise what
you extract from your samples, and \emph{never revisiting a sample is an option}. We believe one can
show theoretically that this \textbf{eliminates one of the classical errors} in the statistical-
learning error decomposition --- the estimation error tied to the finiteness of the dataset.

\begin{figure}[H]
  \centering
  \includegraphics[width=0.80\linewidth]{poisson/data_scaling/scaling_test_rl2.png}
  \caption{Test relative $L^2$ vs.\ training-set size at a fixed optimisation budget. Error decreases
  with dataset size and saturates at the infinite-data (streaming, never-revisit) limit, the dashed
  black line.}
  \label{fig:data-scaling}
\end{figure}

\section{Allen--Cahn}

We now move to the time-dependent setting: an FNO trained as a one-step Allen--Cahn integrator and
rolled out autoregressively.

\paragraph{Opening question --- backprop through time.} Do we know the best way to handle
backpropagation through the unrolled rollout? Here we \textbf{detached the inputs and never
backpropagated through the unrolled operator} --- the \emph{pushforward trick} (Brandstetter et
al.). It performed best across all losses in our tests, and it is cheaper. This deserves discussion,
but it is our current default.

\paragraph{Setup.} We train on a \textbf{10-step} rollout, $128^2$ grid, $\mathrm{dt}=0.01$, and
reaction strength $\varepsilon = 32$ --- note $\varepsilon$ enters the equation \emph{squared}, and
here it is large, not small. Data generation now defaults to a \textbf{convex--concave (Eyre)
splitting} on the $128^2$ grid at this time step, and we train the FNO in exactly the same setting.
From that splitting we derive two physics-informed losses: the natural \textbf{minimizing-movement}
objective, and its \textbf{least-squares} variant. We abandoned the implicit-Euler scheme. Finally
we benchmark against \textbf{data-driven} training. The least-squares variant is preconditioned with
multigrid again --- the problem is nonlinear, but we hand-wave the Jacobian into a constant operator
so multigrid applies.

\paragraph{Findings.}
\begin{itemize}
  \item Physics-informed \textbf{least squares without preconditioning fails} (those curves are not
        in the figure yet).
  \item The \textbf{other three all work quite well} (Figure~\ref{fig:ac-rollout}). Minimizing
        movement appears to have an edge, though that may be down to learning-rate details.
  \item \textbf{Rollouts stay stable out to 100 steps} even though the models were trained on only
        10 steps --- which is quite nice.
\end{itemize}

\begin{figure}[H]
  \centering
  \includegraphics[width=0.95\linewidth]{allen_cahn/realistic_ac/rollout_error.png}
  \caption{Allen--Cahn, relative $L^2$ vs.\ the convex--concave reference along the rollout (solid =
  test, dashed = train), one colour per training loss. Trained on 10 steps, rolled out to 100; the
  error stays controlled well past the training horizon.}
  \label{fig:ac-rollout}
\end{figure}

For a qualitative feel, Figure~\ref{fig:ac-heatmap} shows the 2D fields: rows are rollout timesteps,
columns are the reference next to each learned stepper.

\begin{figure}[H]
  \centering
  \includegraphics[width=0.5\linewidth]{allen_cahn/realistic_ac/heatmaps/heatmap_test_sample0.png}
  \caption{Allen--Cahn rollout fields for one test sample: rows are timesteps, columns are the
  convex--concave reference and each trained model, annotated with the per-panel relative $L^2$.}
  \label{fig:ac-heatmap}
\end{figure}

\section{Future plans}

\paragraph{Stokes.} First results using a \textbf{block-diagonal preconditioner} are promising. We
expect this to land next week.

\paragraph{Fine-tune Poseidon} We can finetune Poseidon physics informed for some PDEs, certainly the Poisson.

\end{document}
